\textbf{SKYFILMS} dispone de los equipos, herramientas e insumos necesarios para realizar de manera segura, precisa y eficiente todas las actividades de mantenimiento preventivo, correctivo y de inspección de sus \textbf{Sistemas de Aeronaves No Tripuladas (UAS)}, el uso de herramientas y equipos adecuados garantiza la integridad técnica de los UAS, la seguridad del personal de mantenimiento.

\section{Tipos de Equipos y Herramientas}
Los equipos y herramientas utilizados por SKYFILMS se clasifican en las siguientes categorías:
\begin{itemize}
    \item \textbf{Herramientas generales:} destornilladores, llaves hexagonales, pinzas, alicates, llaves de torque, cortadores, entre otros.

    \item \textbf{Herramientas y equipos especializados:} medidores de voltaje y corriente, multímetros, y herramientas específicas de los fabricantes de UAS.

    \item \textbf{Equipos de carga y monitoreo de baterías:} cargadores balanceados, estaciones de carga controladas y sistemas de monitoreo de temperatura.

    \item \textbf{Elementos de protección y sujeción:} soportes para mantenimiento, alfombrillas antiestáticas, guantes dieléctricos, y superficies aisladas de trabajo.
\end{itemize}

Cada herramienta o equipo cuenta con una identificación única, asignada en el Registro de Control de Herramientas, que permite su trazabilidad y control operativo.\\

\section{Control y Mantenimiento de Herramientas}
El \textbf{Técnico de UAS} son responsables del control, custodia y mantenimiento de todas las herramientas y equipos utilizados.

Para ello, SKYFILMS aplica las siguientes directrices:

\begin{itemize}
    \item Cada herramienta es registrada en el inventario técnico con código, descripción, condición, responsable y ubicación.

    \item Antes y después de cada jornada de trabajo, el personal técnico debe realizar una verificación visual y funcional del estado de las herramientas.

    \item Las herramientas que presenten daños, desgaste o mal funcionamiento se retiran \textbf{inmediatamente} de servicio y se registran en el Registro de Herramientas y Equipos (F-019), en la columna de estado \textit{No conforme}.

    \item El Técnico UAS es responsable de verificar la disponibilidad mínima de herramientas requeridas por los fabricantes de UAS para los procedimientos de mantenimiento.
\end{itemize}

\textbf{NOTA:} El listado de herramientas se controla mediante el \textbf{Anexo 3 - Registro de Herramientas.}




\section{Selección de Herramientas por Tarea y Modelo}
Antes de cada trabajo, el Técnico UAS elige las herramientas según la tarea y el modelo de aeronave, siguiendo lo que indique el manual de usuario o la guía del fabricante. Ninguna herramienta distinta de las de esta tabla se usa sobre la aeronave sin que el Técnico UAS la registre antes en el F-019.

\begin{longtable}{|>{\raggedright\arraybackslash}p{3.2cm}|>{\raggedright\arraybackslash}p{4.6cm}|>{\raggedright\arraybackslash}p{5.9cm}|}
\hline
\rowcolor[gray]{0.9}\textbf{Tarea} & \textbf{Mavic 4 Pro Creator Combo} & \textbf{BetaFPV Pavo20 Pro II} \\ \hline
\endhead
Cambio de hélices & Hélices originales DJI y el destornillador que indique el manual & Hélices compatibles con el Pavo20 y llave o destornillador hexagonal del tamaño que indique BetaFPV \\ \hline
Inspección de motores, brazos y estructura & Linterna y lupa; sin desarmar & Linterna, lupa y multímetro para medir la resistencia de las bobinas de los motores \\ \hline
Baterías & Cargador y baterías originales DJI; revisión en la aplicación DJI Fly & Cargador balanceado, medidor de celdas, bolsa ignífuga para carga y transporte \\ \hline
Limpieza & Brocha antiestática, paño de microfibra y aire a baja presión; sin líquidos sobre sensores & Brocha antiestática, paño de microfibra, alcohol isopropílico en poca cantidad; limpiar sobre el tapete antiestático \\ \hline
Firmware y configuración & Computador o control DJI RC Pro 2 con DJI Fly o DJI Assistant & Computador con Betaflight Configurator y ExpressLRS Configurator; solo versiones registradas en el capítulo de software del Manual de Operaciones \\ \hline
Soldadura y desensamble profundo & No se realiza en Skyfilms & No se realiza en Skyfilms mientras no exista una organización de mantenimiento aceptada \\ \hline
\end{longtable}

\section{Inventario, Identificación, Calibración y Almacenamiento}
\begin{itemize}
    \item \textbf{Inventario:} el F-019 lista cada herramienta y equipo con código, descripción, modelo de aeronave para el que sirve, estado, responsable y lugar de guardado. El kit de entrega y reintegro del dron es un control distinto y no reemplaza este inventario.
    \item \textbf{Identificación:} cada herramienta lleva una etiqueta o marca con su código. Los equipos de medida (multímetro, medidor de celdas y cualquier torquímetro) llevan además la fecha de la última verificación.
    \item \textbf{Calibración:} cuando el fabricante del instrumento la indique, se sigue su plazo. Si no la indica, el Técnico UAS verifica cada año el equipo contra un valor conocido (por ejemplo, una resistencia o una batería de referencia) y anota el resultado en el F-019. Un equipo vencido o que no pase la verificación no se usa para aceptar o rechazar componentes.
    \item \textbf{Almacenamiento y protección:} las herramientas se guardan en el taller de Ibagué, en estuche o caja cerrada y seca, lejos de las baterías. Los instrumentos de medida se guardan en su estuche. Las piezas pequeñas (tornillos) van en recipientes rotulados por modelo.
    \item \textbf{Instrucciones del fabricante:} el Técnico UAS consulta el manual de usuario del modelo antes de la tarea y no usa una herramienta o un par de apriete diferente al indicado.
\end{itemize}

\section{Insumos Técnicos y Materiales}
Los insumos y materiales de mantenimiento empleados por SKYFILMS —como lubricantes, limpiadores, adhesivos, repuestos, y componentes electrónicos— cumplen con las especificaciones técnicas establecidas por los fabricantes de UAS y son adquiridos únicamente a través de proveedores aprobados.\\

Las políticas de uso incluyen:

\begin{itemize}
    \item Verificación previa de fecha de vencimiento, lote y compatibilidad del producto con el modelo de aeronave.

    \item Almacenamiento controlado, en áreas limpias, secas y ventiladas, conforme a las recomendaciones de seguridad y hojas de datos del fabricante.

    \item Eliminación o disposición final de insumos caducados bajo normas ambientales vigentes.
\end{itemize}

\section{Responsabilidades}

\begin{itemize}
    \item El \textbf{Técnico de UAS} es responsable del cumplimiento general de este capítulo y de la implementación de los controles técnicos sobre equipos, herramientas y administra los inventarios coordinando la reposición o reparación de herramientas.
\end{itemize}

\section{Disposiciones Finales}

Toda herramienta, equipo o insumo técnico que no cumpla con las especificaciones del fabricante o con los estándares de seguridad definidos en este manual deberá ser inmediatamente retirado de servicio.\\

El incumplimiento de las disposiciones de este capítulo será considerado una no conformidad técnica y estará sujeto al proceso interno de \textbf{acciones correctivas y disciplinarias} según los procedimientos de calidad y seguridad operacional de \textbf{SKYFILMS}.
